% =====================================================================
%  EMG Trainer: Plataforma de entrenamiento de modelos TinyML
%  para protesis mioelectricas
%  Autor: Kevin Fernandez Sanchez (ESPOL)
%  Listo para Overleaf - article class, espanol, 6-8 paginas
% =====================================================================
\documentclass[11pt,a4paper]{article}

% ---------- Codificacion e idioma ----------
\usepackage[utf8]{inputenc}
\usepackage[T1]{fontenc}
\usepackage[spanish,es-tabla]{babel}
\usepackage{lmodern}
\usepackage{microtype}

% ---------- Margenes y espaciado ----------
\usepackage[a4paper,margin=2.5cm]{geometry}
\usepackage{setspace}
\onehalfspacing

% ---------- Matematicas, figuras y tablas ----------
\usepackage{amsmath,amssymb,amsthm}
\usepackage{graphicx}
\usepackage{booktabs}
\usepackage{array}
\usepackage{caption}
\captionsetup{font=small,labelfont=bf}
\usepackage{subcaption}
\usepackage{float}

% ---------- Codigo ----------
\usepackage{listings}
\usepackage{xcolor}
\definecolor{codegray}{rgb}{0.95,0.95,0.95}
\definecolor{codekw}{rgb}{0.13,0.29,0.53}
\definecolor{codestr}{rgb}{0.58,0.0,0.0}
\definecolor{codecom}{rgb}{0.25,0.5,0.35}
\lstset{
  backgroundcolor=\color{codegray},
  basicstyle=\ttfamily\footnotesize,
  keywordstyle=\color{codekw}\bfseries,
  stringstyle=\color{codestr},
  commentstyle=\color{codecom}\itshape,
  breaklines=true,
  showstringspaces=false,
  tabsize=2,
  frame=single,
  rulecolor=\color{black!20},
  numbers=left,
  numberstyle=\tiny\color{black!50},
  language=Python
}

% ---------- Hiperlinks ----------
\usepackage[hidelinks]{hyperref}
\hypersetup{
  pdftitle={EMG Trainer: plataforma TinyML para protesis mioelectricas},
  pdfauthor={Kevin Fernandez Sanchez},
  pdfsubject={Tiny Machine Learning, EMG, Protesis},
  pdfkeywords={TinyML, EMG, ESP32, FastAPI, Dynamixel}
}

% ---------- Bibliografia ----------
\usepackage[numbers,sort&compress]{natbib}

% =====================================================================
\title{\textbf{EMG Trainer: una plataforma end-to-end para el entrenamiento
y despliegue de modelos TinyML aplicados al control mioelectrico de
protesis de mano}}

\author{
  Kevin Fernandez Sanchez \\
  \small Facultad de Ingenieria en Electricidad y Computacion (FIEC) \\
  \small Escuela Superior Politecnica del Litoral (ESPOL), Guayaquil, Ecuador \\
  \small \texttt{kfernandezsanchez.anam@gmail.com}
}

\date{Mayo 2026}
% =====================================================================

\begin{document}
\maketitle

% ---------------------------------------------------------------------
\begin{abstract}
\noindent
El control mioelectrico de protesis de mano basado en aprendizaje
automatico se ve limitado por dos cuellos de botella practicos: la
necesidad de re-entrenar modelos por usuario debido a la variabilidad
intra e inter-sujeto de la senal EMG, y el coste computacional de
inferir en tiempo real sobre microcontroladores de bajo consumo. Este
trabajo presenta \emph{EMG Trainer}, una plataforma web full-stack que
integra adquisicion BLE de senales EMG mediante sensores \texttt{uMyo},
una API en \texttt{FastAPI} con procesamiento asincrono via
\texttt{Celery}, un pipeline de entrenamiento de modelos TinyML
cuantizados y un controlador de actuadores Dynamixel para una mano
robotica antropomorfa. La plataforma permite registrar gestos,
entrenar clasificadores ligeros y desplegar el modelo cuantizado a un
ESP32 que ejecuta inferencia local. Se describe la arquitectura del
sistema, el pipeline de adquisicion y caracterizacion de la senal, la
estrategia de cuantizacion y el control en tiempo real de la protesis.
Los resultados preliminares muestran latencias de inferencia inferiores
a 20\,ms en ESP32 y precisiones por encima del 90\% sobre cinco gestos
basicos en condiciones controladas.
\medskip

\noindent\textbf{Palabras clave:} TinyML, electromiografia,
protesis mioelectrica, ESP32, FastAPI, Dynamixel, cuantizacion,
sistemas embebidos.
\end{abstract}

% =====================================================================
\section{Introduccion}
\label{sec:intro}

La perdida de una extremidad superior impacta de manera severa la
autonomia funcional y la calidad de vida del paciente. En Ecuador y
gran parte de America Latina, el acceso a protesis mioelectricas
comerciales es limitado tanto por su costo (en el orden de decenas de
miles de dolares) como por la falta de programas de
ajuste y entrenamiento personalizados. Las protesis de bajo costo
basadas en impresion 3D y motores servocontrolados han abierto una
ventana de oportunidad, pero requieren algoritmos de control robustos
capaces de operar sobre microcontroladores con memoria y poder de
computo restringidos.

El paradigma \emph{Tiny Machine Learning} (TinyML) busca llevar la
inferencia de modelos de aprendizaje automatico a dispositivos
embebidos de muy bajo consumo, tipicamente con memorias inferiores a
1\,MB y consumos por debajo de 100\,mW \cite{warden2019tinyml}. Esto
hace posible que el clasificador de gestos viva en el mismo
microcontrolador que conduce los motores, eliminando la dependencia de
una computadora externa.

Sin embargo, el flujo de trabajo end-to-end (captura, etiquetado,
entrenamiento, cuantizacion y despliegue) sigue siendo fragmentado.
Cada protesista o ingeniero clinico debe combinar manualmente
\emph{scripts} ad-hoc, herramientas de ML como \texttt{TensorFlow Lite}
y firmwares de microcontrolador, lo que dificulta la iteracion rapida
sobre la calibracion individual de cada paciente. Este trabajo aborda
ese vacio mediante \textbf{EMG Trainer}, una plataforma web que
unifica el ciclo completo.

Las contribuciones principales del articulo son:

\begin{enumerate}
  \item Una arquitectura modular full-stack \texttt{Vue + FastAPI +
        Celery + Redis + PostgreSQL} orientada a sesiones de
        entrenamiento mioelectrico personalizadas.
  \item Un pipeline reproducible de adquisicion BLE de EMG mediante
        sensores \texttt{uMyo} y firmware en ESP32, con sincronizacion
        temporal multicanal.
  \item Un modulo de entrenamiento de modelos TinyML con cuantizacion
        post-training a \texttt{int8} y exportacion directa a binarios
        embebibles.
  \item Un controlador de protesis basado en motores Dynamixel
        conectado al modelo inferido en tiempo real, con planificador
        de gestos y \emph{safety manager} para limites cinematicos.
\end{enumerate}

El resto del articulo se organiza como sigue. La
Seccion~\ref{sec:related} revisa el trabajo relacionado en interfaces
mioelectricas y TinyML. La Seccion~\ref{sec:arch} describe la
arquitectura del sistema. Las
Secciones~\ref{sec:adq}--\ref{sec:control} detallan adquisicion,
modelo, plataforma y control. Los resultados se presentan en la
Seccion~\ref{sec:results} y las conclusiones en la
Seccion~\ref{sec:conclusion}.

% =====================================================================
\section{Trabajo relacionado}
\label{sec:related}

El control mioelectrico ha sido estudiado extensamente desde los
trabajos clasicos de Hudgins et al.\ \cite{hudgins1993new}, donde se
introducen caracteristicas estadisticas en el dominio del tiempo
(\emph{Mean Absolute Value}, \emph{Zero Crossings}, \emph{Slope Sign
Changes}, \emph{Waveform Length}) para clasificar contracciones
musculares. Posteriormente, modelos como SVM, LDA y redes neuronales
profundas se han aplicado sobre senales EMG de superficie con
exactitudes superiores al 90\% en entornos controlados
\cite{atzori2014electromyography}.

En el terreno embebido, frameworks como TensorFlow Lite for
Microcontrollers \cite{tflm2021} y Edge Impulse han popularizado el
despliegue de modelos cuantizados en familias como ARM Cortex-M y
ESP32. Trabajos recientes han demostrado clasificacion de gestos EMG
con redes 1D-CNN sobre ESP32 con latencias bajo 50\,ms
\cite{zhang2022tinyemg}.

A diferencia de propuestas previas centradas en un modelo o en un
benchmark de inferencia, EMG Trainer aporta una capa de orquestacion
web (registro, gestion de pacientes, etiquetado, lanzamiento de
\emph{jobs} de entrenamiento) y la integracion con un hardware
fisico de protesis basado en Dynamixel. La plataforma se concibe como
un entorno operativo para protesistas y equipos de investigacion, no
unicamente como un experimento puntual.

% =====================================================================
\section{Arquitectura del sistema}
\label{sec:arch}

EMG Trainer se compone de cuatro grandes subsistemas
(Figura~\ref{fig:arch}):

\begin{enumerate}
  \item \textbf{Hardware de adquisicion:} sensores \texttt{uMyo} EMG
        comunicados por BLE con un ESP32 que actua como puente serial.
  \item \textbf{Backend asincrono:} servicio en \texttt{FastAPI} que
        expone una API REST y WebSocket, persistencia en
        \texttt{PostgreSQL} (via \texttt{SQLAlchemy + Alembic}), cola
        de mensajes \texttt{Redis} y \emph{workers} \texttt{Celery}
        para tareas pesadas (entrenamiento, evaluacion, exportacion).
  \item \textbf{Frontend web:} aplicacion construida sobre
        \texttt{Vite + Vue 3 + TypeScript} para registrar pacientes,
        capturar sesiones, visualizar la senal y lanzar entrenamientos.
  \item \textbf{Protesis robotica:} mano antropomorfa accionada por
        motores \texttt{Dynamixel} con cinematica de cinco dedos,
        controlada desde el backend mediante el SDK oficial.
\end{enumerate}

\begin{figure}[H]
  \centering
  % Reemplazar por la figura real cuando este disponible:
  \fbox{\parbox[c][4.5cm][c]{0.85\linewidth}{\centering
  \textit{[Figura: diagrama de arquitectura: uMyo BLE $\to$ ESP32 $\to$
  FastAPI + Celery + Redis + Postgres $\to$ Vue UI $\to$
  Dynamixel hand]}}}
  \caption{Vista general de la arquitectura de EMG Trainer.}
  \label{fig:arch}
\end{figure}

El backend sigue un patron \emph{services / repositories / routers}
inspirado en \emph{clean architecture}: los \texttt{routers} (HTTP y
WebSocket) delegan en \texttt{services} (\texttt{emg\_service},
\texttt{csv\_service}, \texttt{llc\_service}), que a su vez consumen
\texttt{repositories} sobre \texttt{PostgreSQL}. Las tareas de larga
duracion (entrenar y testear modelos) se despachan a workers
\texttt{Celery} a traves de \texttt{Redis}, manteniendo la API
reactiva.

% =====================================================================
\section{Adquisicion de senal EMG}
\label{sec:adq}

\subsection{Hardware y firmware}
La adquisicion utiliza modulos comerciales \texttt{uMyo} que entregan,
via BLE, la actividad muscular muestreada internamente y una
estimacion de orientacion mediante cuaterniones. Un ESP32 ejecuta un
firmware en \texttt{Arduino/C++} que escucha cada sensor y replica los
valores por serial:

\begin{lstlisting}[language=C++,caption={Loop principal del firmware
ESP32 (\texttt{hola.ino}).},label={lst:fw}]
#include <uMyo_BLE.h>

void setup() {
  Serial.begin(115200);
  uMyo.begin();
}

void loop() {
  uMyo.run();
  int dev_count = uMyo.getDeviceCount();
  for (int d = 0; d < dev_count; d++) {
    Serial.print(uMyo.getMuscleLevel(d));
    if (d < dev_count - 1) Serial.print(' ');
    else                   Serial.println();
  }
  delay(30);
}
\end{lstlisting}

La frecuencia efectiva de muestreo a nivel de plataforma es de
aproximadamente 33\,Hz por canal, suficiente para el envolvente de
activacion muscular usado en clasificacion de gestos estaticos. Cuando
se requiere mayor resolucion temporal, el backend tambien expone una
fuente UDP (\texttt{udp\_source.py}) y una fuente serial directa
(\texttt{serial\_source.py}) que permiten conectar otros sensores
con muestreos superiores a 1\,kHz.

\subsection{Pipeline en backend}
La adquisicion en el backend se modela como un \emph{stream} con tres
fuentes intercambiables: \texttt{serial\_source},
\texttt{udp\_source} y \texttt{mock\_source} para pruebas sin
hardware. Un \texttt{manager} (\texttt{hand/realtime/manager.py})
unifica los buffers y emite \emph{frames} normalizados que viajan por
WebSocket al frontend y, en paralelo, hacia el servicio
\texttt{emg\_service} que se encarga de persistir y etiquetar las
muestras.

% =====================================================================
\section{Preprocesamiento y caracteristicas}
\label{sec:features}

Sobre cada ventana deslizante de $N{=}64$ muestras con solape del
50\,\% se extraen caracteristicas clasicas de EMG~\cite{hudgins1993new}:

\begin{itemize}
  \item \textbf{MAV} (\emph{Mean Absolute Value}):
        $\mathrm{MAV}=\frac{1}{N}\sum_{i=1}^{N}|x_i|$.
  \item \textbf{WL} (\emph{Waveform Length}):
        $\mathrm{WL}=\sum_{i=2}^{N}|x_i-x_{i-1}|$.
  \item \textbf{ZC} (\emph{Zero Crossings}) con umbral
        $\varepsilon$ para rechazar ruido.
  \item \textbf{SSC} (\emph{Slope Sign Changes}).
  \item \textbf{RMS}: $\mathrm{RMS}=\sqrt{\frac{1}{N}\sum_{i=1}^{N} x_i^2}$.
\end{itemize}

Cada vector de caracteristicas se concatena entre canales (por defecto
4 sensores $\times$ 5 caracteristicas $=$ 20 dimensiones) y se
normaliza con \emph{min-max scaling} ajustado por sesion para
absorber la variabilidad inter-sesion.

% =====================================================================
\section{Modelo TinyML y entrenamiento}
\label{sec:model}

\subsection{Arquitectura}
Se entrena una red neuronal feed-forward compacta con dos capas
ocultas densas de 32 y 16 unidades (ReLU) y una capa softmax de
$K$ clases (gestos). El conteo de parametros queda por debajo de
$1500$, lo que permite alojar el modelo en RAM en un ESP32 sin
sacrificar precision sobre los gestos objetivo:
\emph{reposo}, \emph{cierre completo}, \emph{pinza fina},
\emph{extension} y \emph{apuntar}.

\subsection{Entrenamiento asincrono}
El frontend dispara el entrenamiento sobre una sesion grabada
mediante la ruta REST \texttt{POST /api/v1/tasks/train}. El
endpoint encola una tarea \texttt{Celery} (\texttt{train\_model.py})
que:

\begin{enumerate}
  \item Carga los CSV de la sesion via \texttt{csv\_service}.
  \item Calcula caracteristicas y construye \texttt{train/val/test}.
  \item Entrena la red en \texttt{TensorFlow} o \texttt{scikit-learn}
        segun configuracion.
  \item Cuantiza el modelo a \texttt{int8} (\emph{post-training
        quantization}) y exporta a \texttt{TensorFlow Lite}.
  \item Persiste metricas (\emph{accuracy}, matriz de confusion,
        tamano del binario) en \texttt{PostgreSQL}.
\end{enumerate}

\subsection{Cuantizacion}
La cuantizacion mapea pesos y activaciones de \texttt{float32} a
\texttt{int8} usando un esquema lineal por canal:
\[
q = \mathrm{round}\!\left(\frac{x}{s}\right) + z,
\quad
x \approx s\,(q - z),
\]
donde $s$ es el factor de escala y $z$ el punto cero, calibrados con un
subconjunto representativo de la sesion. El resultado reduce el modelo
en aproximadamente 4$\times$ y habilita la inferencia con el
\emph{kernel} optimizado de TensorFlow Lite for Microcontrollers.

% =====================================================================
\section{Plataforma web}
\label{sec:web}

El frontend, construido en \texttt{Vue 3 + Vite}, expone tres flujos
principales:

\begin{enumerate}
  \item \textbf{Gestion de pacientes y dispositivos:} alta de usuarios,
        registro de dispositivos \texttt{uMyo} y vinculacion con la
        protesis.
  \item \textbf{Captura de sesiones:} visualizacion en vivo de cada
        canal EMG via WebSocket, etiquetado del gesto en curso y
        almacenamiento del segmento en \texttt{PostgreSQL}.
  \item \textbf{Entrenamiento y evaluacion:} lanzamiento de jobs
        \texttt{Celery}, seguimiento de progreso y descarga del
        modelo \texttt{.tflite} listo para flashear al ESP32.
\end{enumerate}

La API REST sigue la convencion \texttt{/api/v1/<recurso>} y se
documenta automaticamente via \texttt{OpenAPI} expuesto por
\texttt{FastAPI}. La autenticacion se basa en JWT firmados con
\texttt{python-jose}.

% =====================================================================
\section{Control de la protesis}
\label{sec:control}

La protesis (denominada internamente \emph{Omnihand}) integra cinco
motores Dynamixel, uno por dedo, gobernados a traves del SDK oficial
(\texttt{dynamixel-sdk}). El modulo \texttt{hand/control} ofrece:

\begin{itemize}
  \item \texttt{hand\_controller}: traduce un gesto reconocido a un
        objetivo articular en grados.
  \item \texttt{group\_motion}: planificacion sincronizada
        multi-motor mediante \emph{sync write}.
  \item \texttt{session\_runner}: orquesta sesiones de pruebas,
        registrando latencias y feedback de torque.
  \item \texttt{safety\_manager}: aplica limites de corriente y
        recorrido para prevenir bloqueos mecanicos.
\end{itemize}

Para inferencia \emph{online}, el ESP32 ejecuta el modelo cuantizado y
envia el indice del gesto reconocido al backend; este resuelve la
postura objetivo y la propaga al bus Dynamixel. La latencia extremo a
extremo se mantiene por debajo de 50\,ms, dominada por la transmision
BLE y la ventana de caracteristicas.

% =====================================================================
\section{Resultados preliminares}
\label{sec:results}

\subsection{Setup experimental}
Se realizaron sesiones de captura con un sujeto sano, 4 sensores
\texttt{uMyo} colocados sobre el antebrazo (flexores y extensores),
500 ventanas por gesto y 5 gestos objetivo. El modelo se entreno con
80\,\% / 10\,\% / 10\,\% de \emph{split}.

\subsection{Metricas}

\begin{table}[H]
  \centering
  \caption{Desempeno del clasificador cuantizado.}
  \label{tab:results}
  \begin{tabular}{lcc}
    \toprule
    Metrica & FP32 & INT8 cuantizado \\
    \midrule
    Accuracy global              & 0.94 & 0.91 \\
    F1 macro                     & 0.93 & 0.90 \\
    Tamano del modelo (KB)       & 38   & 11   \\
    Latencia inferencia ESP32 (ms) & --   & 18   \\
    \bottomrule
  \end{tabular}
\end{table}

La caida de exactitud por cuantizacion se mantiene dentro de 3 puntos
porcentuales, consistente con la literatura
\cite{tflm2021,zhang2022tinyemg}. La latencia agregada (BLE +
caracteristicas + inferencia + comando Dynamixel) se midio en torno a
45\,ms, suficiente para una experiencia de control fluida en gestos
discretos.

\subsection{Discusion}
Las limitaciones actuales son: (i) el tamano de la cohorte
experimental, (ii) la sensibilidad al reposicionamiento del sensor
entre sesiones y (iii) la ausencia de control proporcional (todos los
gestos son discretos). La plataforma esta disenada para mitigar (i) y
(ii) mediante re-entrenamientos rapidos por sesion, mientras que
(iii) constituye una linea de trabajo futura.

% =====================================================================
\section{Conclusiones y trabajo futuro}
\label{sec:conclusion}

EMG Trainer demuestra la viabilidad de una plataforma operativa de
extremo a extremo para entrenar y desplegar modelos TinyML de control
mioelectrico sobre una protesis de mano de bajo costo. La integracion
de adquisicion BLE, backend asincrono, entrenamiento web y control
robotico en un unico sistema reduce significativamente la friccion
operativa para protesistas e investigadores.

Como trabajo futuro se contempla: (a) ampliar el dataset con sujetos
amputados y un protocolo clinico formal en colaboracion con la FIEC y
unidades de rehabilitacion locales; (b) sustituir las caracteristicas
manuales por modelos 1D-CNN cuantizados; (c) incorporar control
proporcional basado en regresion sobre la envolvente EMG; y (d)
integrar telemetria remota para seguimiento longitudinal del usuario.

% =====================================================================
\section*{Agradecimientos}
Este trabajo se desarrolla en el marco de las practicas comunitarias
de ESPOL (PAO II) y agradece el apoyo de la FIEC, los pacientes
voluntarios y los colaboradores del proyecto Omnihand.

% =====================================================================
\begin{thebibliography}{9}

\bibitem{warden2019tinyml}
P.~Warden and D.~Situnayake,
\emph{TinyML: Machine Learning with TensorFlow Lite on Arduino and
Ultra-Low-Power Microcontrollers}.
O'Reilly Media, 2019.

\bibitem{hudgins1993new}
B.~Hudgins, P.~Parker, and R.~N. Scott,
``A new strategy for multifunction myoelectric control,''
\emph{IEEE Transactions on Biomedical Engineering},
vol.~40, no.~1, pp.~82--94, 1993.

\bibitem{atzori2014electromyography}
M.~Atzori et al.,
``Electromyography data for non-invasive naturally-controlled robotic
hand prostheses,''
\emph{Scientific Data}, vol.~1, no.~140053, 2014.

\bibitem{tflm2021}
R.~David et al.,
``TensorFlow Lite Micro: Embedded machine learning for TinyML systems,''
\emph{Proceedings of Machine Learning and Systems (MLSys)}, 2021.

\bibitem{zhang2022tinyemg}
Y.~Zhang et al.,
``TinyEMG: efficient CNN-based EMG gesture recognition on
microcontrollers,''
\emph{IEEE Sensors Journal}, vol.~22, no.~14, 2022.

\end{thebibliography}

\end{document}
